% Part: sets-functions-relations
% Chapter: infinite
% Section: dedekind
%
\documentclass[../../../include/open-logic-section]{subfiles}

\begin{document}

\olfileid{sfr}{infinite}{dedekind}
\olsection{Dedekind 代数}

我们不仅希望自然数是无限的；还希望它们具有特定的（代数）性质：它们需要在加法、乘法等运算下表现良好。

Dedekind 的想法是将\emph{后继函数}的概念视为基础，然后用以下性质来刻画这些数：
\begin{enumerate}
	\item 存在一个数 $0$，它不是任何数的后继，
	\\即 $0 \notin \ran{s}$
	\\即 $\forall x\ s(x) \neq 0$
	\item 不同的数有不同的后继，
	\\即 $s$ 是单射
	\\即 $\forall x \forall y (s(x) = s(y) \lif x = y)$
	\item\ollabel{repeatedapplication} 每个数都是通过对 $0$ 反复应用后继函数得到的。
\end{enumerate}
前两个条件用一阶逻辑很容易处理（见上文）。但我们无法仅用一阶逻辑处理 \olref{repeatedapplication}。Dedekind 的突破在于以集合论的方式将条件 \olref{repeatedapplication} 重新表述如下：
\begin{enumerate}
	\item[3$'$.] 自然数是最小的\emph{对后继函数封闭的}集合：也就是说，如果我们对集合中的任何元素应用 $s$，就会得到该集合中的另一个元素，
\end{enumerate}

但我们需要仔细阐明这一点。

\begin{defn}\ollabel{Closure}
	对于任意函数 $f$，集合 $X$ 是 $f$-\emph{封闭的}，当且仅当
	$(\forall x \in X)f(x) \in X$。现对任意 $o$ 定义：
	$$\closureofunder{f}{o} = \bigcap\Setabs{X}{o \in X\text{ and }X
\text{ is $f$-closed}}$$
\end{defn}


所以 $\closureofunder{f}{o}$ 就是所有以 $o$ 为元素的 $f$-封闭集合的交集。因此直观上，
$\closureofunder{f}{o}$ 就是以 $o$ 为元素的\emph{最小} $f$-封闭集合。下面的结果将这一直观想法精确化：
\begin{lem}\ollabel{closureproperties}
	对于任意函数 $f$ 和任意 $o \in A$：
	\begin{enumerate}
		\item\ollabel{closurehaselem} $o \in \closureofunder{f}{o}$；且
		\item\ollabel{closureclosed} $\closureofunder{f}{o}$ 是 $f$-封闭的；且
		\item\ollabel{closuresmallest} 若 $X$ 是 $f$-封闭的且 $o \in
		X$，则 $\closureofunder{f}{o} \subseteq X$。
	\end{enumerate}
\end{lem}

\begin{proof}
注意到至少存在一个以 $o$ 为元素的 $f$-封闭集合，即 $\ran{f}\cup
\{o\}$。所以 $\closureofunder{f}{o}$，即\emph{所有}此类集合的交集，是存在的。
现在我们需要逐一验证
\olref{closurehaselem}--\olref{closuresmallest}。

关于 \olref{closurehaselem}：$o \in \closureofunder{f}{o}$，因为它是若干均包含 $o$ 的集合的交集。

关于 \olref{closureclosed}：假设 $x \in \closureofunder{f}{o}$。那么对于任意满足 $o \in X$ 且 $X$ 是 $f$-封闭的 $X$，都有 $x \in X$；而由于 $X$ 是
$f$-封闭的，从而 $f(x) \in X$。所以 $f(x) \in \closureofunder{f}{o}$。

关于 \olref{closuresmallest}：一般地，若 $X
\in C$，则 $\bigcap C \subseteq X$。
\end{proof}

利用上述结果，我们可以说：

\begin{defn}
一个 \emph{Dedekind 代数} 是一个集合 $A$，连同一个函数 $f
\colon A \to A$ 和某个 $o \in A$，满足：
	\begin{enumerate}
		\item \ollabel{ded:proper} $o \notin \ran{f}$
		\item \ollabel{ded:injection} $f$ 是单射
		\item \ollabel{ded:closure} $A = \closureofunder{f}{o}$
	\end{enumerate}
\end{defn}

由于 $A = \closureofunder{f}{o}$，我们前面的结果告诉我们，
$A$ 是包含 $o$ 的最小 $f$-封闭集合。显然，一个
Dedekind 代数是 Dedekind 无限的；只需看定义中的条目
\olref{ded:proper} 和 \olref{ded:injection}。但
更令人兴奋的事实是，任何 Dedekind 无限集合都可以转化为一个 Dedekind 代数。

\begin{thm}\ollabel{thm:DedekindInfiniteAlgebra}
如果存在 Dedekind 无限集合，那么就存在 Dedekind 代数。
\end{thm}

\begin{proof}
设 $D$ 是 Dedekind 无限的。于是存在一个单射 $g \colon D \to
D$ 和一个元素 $o  \in D \setminus \ran{g}$。现在令 $A =
\closureofunder{g}{o}$；由 \olref{closureproperties}，$A$ 存在且 $o \in A$。令 $f =
\funrestrictionto{g}{A}$。我们将证明 $A$、$f$、$o$ 构成一个 Dedekind
代数。

关于 \olref{ded:proper}：$o \notin \ran{g}$ 且 $\ran{f}
\subseteq \ran{g}$，因此 $o\notin \ran{f}$。

关于 \olref{ded:injection}：$g$ 是 $D$ 上的单射；所以 $f
\subseteq g$ 必定也是单射。

关于 \olref{ded:closure}：由 \olref{closureproperties}，$A$ 是对 $g$ 封闭的；显然，$A$ 也是对 $f$ 封闭的。因此根据 \olref{closureproperties}，有 $\closureofunder{f}{o} \subseteq A$。由于 $\closureofunder{f}{o}$ 也是对 $f$ 封闭的，且 $f = \funrestrictionto{g}{A}$，可知 $\closureofunder{f}{o}$ 也是对 $g$ 封闭的。于是再由 \olref{closureproperties}，得到 $A \subseteq \closureofunder{f}{o}$。
\end{proof}

\end{document}